% !TEX root = ../main.tex
% ---------------------------------------------------------------
% GENERATED FILE - DO NOT EDIT.
% Produced by docs/tools/render_reference.py from the repository source.
% Regenerate with: make docs
% ---------------------------------------------------------------

\apibreak
\section{\texttt{CircuitBreaker}}
\label{class:feedback_intelligence_agent.resilience.CircuitBreaker}\index{Classes!CircuitBreaker}

In-memory circuit-breaker state for one provider instance.
\apifield{Attributes}
\begin{arglist}
  \item[\texttt{failure\_\allowbreak{}count}] \texttt{int}, default \texttt{0}
  \item[\texttt{state}] \texttt{CircuitState}, default \texttt{'closed'}
  \item[\texttt{opened\_\allowbreak{}at}] \texttt{float \textbar{} None}, default \texttt{None}
\end{arglist}
\apifield{Source}\texttt{src/\allowbreak{}feedback\_\allowbreak{}intelligence\_\allowbreak{}agent/\allowbreak{}resilience.\allowbreak{}py:43}~(\href{https://github.com/DiogoRibeiro7/feedback-intelligence-agent/blob/b6cbc710ffdfa5f26f700c999d594c67588c1e16/src/feedback_intelligence_agent/resilience.py#L43}{online}), dataclass, module \cref{module:feedback_intelligence_agent.resilience}

\subsection{\texttt{before\_call}}
\label{method:feedback_intelligence_agent.resilience.CircuitBreaker.before_call}\index{Methods!before\_call}
Reject calls while open, or move to half-open after recovery.
\apifield{Usage}
\begin{lstlisting}[style=usage, language=Python]
before_call(self, *, now: float, recovery_seconds: float) -> None
\end{lstlisting}
\apifield{Arguments}
\begin{arglist}
  \item[\texttt{now}] \texttt{float}, required, keyword-only.
  \item[\texttt{recovery\_\allowbreak{}seconds}] \texttt{float}, required, keyword-only.
\end{arglist}
\apifield{Raises}\texttt{LLMProviderError}
\apifield{Source}\texttt{src/\allowbreak{}feedback\_\allowbreak{}intelligence\_\allowbreak{}agent/\allowbreak{}resilience.\allowbreak{}py:50}~(\href{https://github.com/DiogoRibeiro7/feedback-intelligence-agent/blob/b6cbc710ffdfa5f26f700c999d594c67588c1e16/src/feedback_intelligence_agent/resilience.py#L50}{online})

\subsection{\texttt{record\_success}}
\label{method:feedback_intelligence_agent.resilience.CircuitBreaker.record_success}\index{Methods!record\_success}
Close the circuit after any successful provider call.
\apifield{Usage}
\begin{lstlisting}[style=usage, language=Python]
record_success(self) -> None
\end{lstlisting}
\apifield{Source}\texttt{src/\allowbreak{}feedback\_\allowbreak{}intelligence\_\allowbreak{}agent/\allowbreak{}resilience.\allowbreak{}py:60}~(\href{https://github.com/DiogoRibeiro7/feedback-intelligence-agent/blob/b6cbc710ffdfa5f26f700c999d594c67588c1e16/src/feedback_intelligence_agent/resilience.py#L60}{online})

\subsection{\texttt{record\_failure}}
\label{method:feedback_intelligence_agent.resilience.CircuitBreaker.record_failure}\index{Methods!record\_failure}
Track a failed call and open the circuit when the threshold is reached.
\apifield{Usage}
\begin{lstlisting}[style=usage, language=Python]
record_failure(self, *, now: float, threshold: int) -> None
\end{lstlisting}
\apifield{Arguments}
\begin{arglist}
  \item[\texttt{now}] \texttt{float}, required, keyword-only.
  \item[\texttt{threshold}] \texttt{int}, required, keyword-only.
\end{arglist}
\apifield{Source}\texttt{src/\allowbreak{}feedback\_\allowbreak{}intelligence\_\allowbreak{}agent/\allowbreak{}resilience.\allowbreak{}py:66}~(\href{https://github.com/DiogoRibeiro7/feedback-intelligence-agent/blob/b6cbc710ffdfa5f26f700c999d594c67588c1e16/src/feedback_intelligence_agent/resilience.py#L66}{online})
